\section{}

\textbf{Tiempo de lectura:} \{s13-reading-time\}

En el \hyperref[memoria_compartida]{???} vimos que varios procesos
pueden compartir regiones de la memoria con el objeto de cooperar en las
tareas que deben desempeñar. Además, en el \hyperref[hilos]{???} vimos
que en los procesos multihilo todos los hilos comparten el espacio de
direcciones del proceso al que pertenecen, lo que significa que pueden
acceder al mismo tiempo a las variables globales y a la memoria
reservada dinámicamente.

Ambas posibilidades introducen algunos riesgos, puesto que el acceso
simultáneo a los datos compartidos puede ocasionar inconsistencias. Así
que ha llegado el momento de discutir cómo se puede asegurar la
ejecución ordenada de hilos o procesos cooperativos que comparten
regiones de la memoria, con el fin de mantener la consistencia de los
datos.

En este capítulo hablaremos de hilos y de procesos que comparten memoria
indistintamente. En ambos casos el problema es el mismo y las soluciones
similares.

\section{El problema de las secciones
críticas}\label{_el_problema_de_las_secciones_cruxedticas}

Llamamos \textbf{{}condición de carrera} a la situación en la que varios
procesos o hilos pueden acceder y manipular los mismos datos al mismo
tiempo ---es decir, de forma \textbf{{}concurrente}--- y donde el
resultado de la ejecución depende del orden particular en el que tienen
lugar dichos accesos. Estas situaciones ocurren frecuentemente en los
sistemas operativos, puesto que diferentes componentes del mismo
manipulan los mismos recursos interfiriendo unos con otros.

\subsection{Problema del
productor-consumidor}\label{_problema_del_productor_consumidor}

Para ilustrarlo, veamos un problema clásico de concurrencia: el
\textbf{problema del {}productor-consumidor}.

Supongamos que dos hilos o procesos comparten una región de la memoria
que contiene un vector de elementos y un contador con el número de
elementos del vector.

El primer hilo realiza varias tareas, que no entraremos a describir. Lo
importante es que, como resultado de esas tareas, en ocasiones añade un
elemento al vector e incrementa el contador que indica el número de
elementos en el vector. Es decir, el primer hilo actúa como un
\textbf{{}productor} de elementos del vector.

A continuación mostramos una porción de la función del productor:

\begin{Shaded}
\begin{Highlighting}[]
\ControlFlowTok{while}\OperatorTok{(}\CommentTok{/* ... */}\OperatorTok{)}
\OperatorTok{\{}
\NormalTok{    item }\OperatorTok{=}\NormalTok{ produce\_item}\OperatorTok{();}

    \CommentTok{// Si el vector está lleno, esperar}
    \ControlFlowTok{while} \OperatorTok{(}\NormalTok{count }\OperatorTok{==}\NormalTok{ VECTOR\_SIZE}\OperatorTok{);}

    \CommentTok{// Añadir el elemento al vector}
\NormalTok{    vector}\OperatorTok{[}\NormalTok{count}\OperatorTok{]} \OperatorTok{=}\NormalTok{ item}\OperatorTok{;}
    \OperatorTok{++}\NormalTok{count}\OperatorTok{;}
\OperatorTok{\}}
\end{Highlighting}
\end{Shaded}

El segundo hilo también realiza varias tareas que no describiremos. Pero
para realizar esas tareas en ocasiones debe tomar un elemento del vector
compartido, decrementando el contador para indicar que ahora hay un
elemento menos en el vector. Es decir, el segundo hilo actúa como un
\textbf{{}consumidor} de elementos del vector.

A continuación mostramos una porción de la función del consumidor:

\begin{Shaded}
\begin{Highlighting}[]
\ControlFlowTok{while}\OperatorTok{(}\CommentTok{/* ... */}\OperatorTok{)}
\OperatorTok{\{}
    \CommentTok{// Si el vector está vacio, esperar}
    \ControlFlowTok{while} \OperatorTok{(}\NormalTok{count }\OperatorTok{==} \DecValTok{0}\OperatorTok{);}

    \CommentTok{// Extraer un elemento del vector}
    \OperatorTok{{-}{-}}\NormalTok{count}\OperatorTok{;}
\NormalTok{    item }\OperatorTok{=}\NormalTok{ vector}\OperatorTok{[}\NormalTok{count}\OperatorTok{];}

\NormalTok{    consume\_item}\OperatorTok{(}\NormalTok{item}\OperatorTok{);}
\OperatorTok{\}}
\end{Highlighting}
\end{Shaded}

Aunque el \textbf{problema del productor-consumidor} parezca artificial,
lo cierto es que es muy común.

Por ejemplo, en una herramienta de grabación de audio, el
\textbf{productor} es un hilo dedicado a obtener bloques de muestras
grabadas a través de la API multimedia del sistema operativo. Mientras
tanto, otro hilo puede dedicarse a tomar las muestras y realizar
diversas transformaciones, como: reducir el ruido, mezclar con otras
fuentes de sonido o aplicar algún tipo de efecto digital. Este segundo
hilo es el \textbf{consumidor}. La manera de conectar ambos es tener un
vector compartido, donde se depositan los bloques de muestras cuando
llegan y de dónde se extraen para su tratamiento. Así, ambos hilos
pueden trabajar a su propio ritmo, de forma casi independiente.

Aunque el código anterior del productor y del consumidor es correcto
cuando no coinciden en el tiempo al ejecutarse, no funciona
adecuadamente cuando sí lo hacen. El motivo es que los dos hilos
comparten y modifican la variable \texttt{count}.

Obviamente, las sentencias \texttt{++count} y \texttt{-\/-count} pueden
modificar \texttt{count} al mismo tiempo si los hilos se ejecutan en
núcleos diferentes en máquinas multiprocesador o multinúcleo. Eso
también puede ocurrir en máquinas monoprocesador con sistemas operativos
con planificación expropiativa (véase el
\hyperref[_planificaciuxf3n_expropiativa]{???}) dado que en esos
sistemas un hilo puede ser interrumpido en cualquier momento para dar
paso a la ejecución de otro.

Por ejemplo, \texttt{++count} podría dividirse por el compilador en las
siguientes operaciones, al generar las instrucciones del procesador:

\textbf{++count}

\begin{verbatim}
registro1 = [count];
registro1 = registro1 + 1;
[count] = registro1;
\end{verbatim}

Donde \texttt{registro1} representa un registro de la CPU. De forma
parecida la sentencia \texttt{-\/-count} podría ser implementada de la
siguiente manera:

\textbf{-\/-count}

\begin{verbatim}
registro2 = [count];
registro2 = registro2 - 1;
[count] = registro2;
\end{verbatim}

Donde nuevamente \texttt{registro2} representa un registro de la CPU.

Aunque \texttt{registro1} y \texttt{registro2} pueden ser el mismo
registro físico, el contenido de los registros se guarda y se recupera
durante los cambios de contexto de un hilo al otro, por lo que es
indiferente si son el mismo o registros diferentes.

En sistemas monoprocesador, el que las sentencias \texttt{++count} y
\texttt{-\/-count} se ejecute de forma concurrente, es similar a que las
instrucciones de lenguaje máquina de ambas sentencias en ambos hilos o
procesos se entrelacen en algún orden aleatorio.

Por ejemplo, un posible entrelazado de las instrucciones en lenguaje
máquina entre hilos, suponiendo que inicialmente \texttt{count\ =\ 5},
podría ser el siguiente:

\begin{Shaded}
\begin{Highlighting}[]
\CommentTok{// Entra ++count}
\NormalTok{registro1 }\OperatorTok{=} \OperatorTok{[}\NormalTok{count}\OperatorTok{];}        \CommentTok{// registro1 = 5}
\NormalTok{registro1 }\OperatorTok{=}\NormalTok{ registro1 }\OperatorTok{+} \DecValTok{1}\OperatorTok{;}  \CommentTok{// registro1 = 6}
\CommentTok{// Sale ++count y entra {-}{-}count}
\NormalTok{registro2 }\OperatorTok{=} \OperatorTok{[}\NormalTok{count}\OperatorTok{];}        \CommentTok{// registro2 = 5}
\NormalTok{registro2 }\OperatorTok{=}\NormalTok{ registro2 }\OperatorTok{{-}} \DecValTok{1}\OperatorTok{;}  \CommentTok{// registro2 = 4}
\CommentTok{// Sale {-}{-}count y entra ++count}
\OperatorTok{[}\NormalTok{count}\OperatorTok{]} \OperatorTok{=}\NormalTok{ registro1}\OperatorTok{;}        \CommentTok{// count = 6 }
\CommentTok{// Entra {-}{-}count}
\OperatorTok{[}\NormalTok{count}\OperatorTok{]} \OperatorTok{=}\NormalTok{ registro2}\OperatorTok{;}        \CommentTok{// count = 4 }
\end{Highlighting}
\end{Shaded}

\begin{itemize}
\item
  Así llegamos al resultado incorrecto \texttt{count\ =\ 4}, indicando
  que hay 4 elementos en el vector cuando realmente hay 5.
\item
  Si invertimos el orden de las sentencias obtendremos el resultado,
  también incorrecto, \texttt{count\ =\ 6}.
\end{itemize}

Como se puede apreciar, hemos llegado a estos valores incorrectos porque
hemos permitido la manipulación concurrente de la variable
\texttt{count}. Según como se entrelacen las instrucciones de
\texttt{++count} y \texttt{-\/-count} en la CPU, el resultado final
podría ser: 4, 5 o 6. Sin embargo, el único resultado correcto es 5, que
es el que obtendremos si ejecutamos las sentencias secuencialmente, sin
mezclar las operaciones en las que se dividen.

\subsection{Manipular estructuras de
datos}\label{_manipular_estructuras_de_datos}

Obviamente, el problema comentado no aparece solo en sentencias simples,
sino también en bloques de código destinados a hacer tareas complejas,
como manipular estructuras de datos.

Por ejemplo, supongamos que \texttt{vector} no es un simple \emph{array}
de elementos, sino una lista enlazada, de tal forma que ahora extraer un
elemento sería así:

\begin{Shaded}
\begin{Highlighting}[]
\OperatorTok{{-}{-}}\NormalTok{count}\OperatorTok{;}
\NormalTok{item }\OperatorTok{=}\NormalTok{ vector}\OperatorTok{.}\NormalTok{extract}\OperatorTok{(}\NormalTok{count}\OperatorTok{);}
\end{Highlighting}
\end{Shaded}

y el método \texttt{extract()} tendría que dar los siguientes pasos:

\begin{enumerate}
\def\labelenumi{\arabic{enumi}.}
\item
  Iterar sobre la lista para buscar el nodo en la posición
  \texttt{count}.
\item
  Al encontrarlo, preservar en variables locales el puntero a ese nodo y
  al previo.
\item
  Cambiar en el nodo previo el puntero al siguiente nodo, para que
  apunte al nodo tras el que queremos extraer. En este momento el nodo a
  extraer ya no pertenece a la lista enlazada.
\item
  Extrae el \texttt{item} del campo que lo contiene en el nodo.
\item
  Destruir el nodo.
\item
  Salir del método retornando el elemento.
\end{enumerate}

Esto genera varios momentos cruciales en torno al paso 3, que puedan dar
lugar a \textbf{condiciones de carrera}. Por ejemplo, si el hilo es
interrumpido tras guardar el puntero al nodo en una variable local y
llega otro hilo que extrae ---y destruye--- antes ese mismo nodo, el
puntero ya no es válido ---es un \emph{dangling pointer} o referencia
colgante--- al continuar la ejecución del primer hilo. Y lo mismo ocurre
con el puntero al nodo previo o al siguiente, si el hilo es interrumpido
y otro hilo destruye antes alguno de ellos.

Los problemas que esto puede causar son diversos, según el momento
exacto en el que ocurra. Puede haberlos al intentar leer el elemento
guardado en el nodo en el paso 4, porque este último ya no exista.
También, al intentar actualizar, en el paso 3, el puntero al siguiente
nodo en el nodo previo, porque el nodo previo no exista. Incluso puede
que la función termine con aparente normalidad, pero dejando que el nodo
previo apunte a un nodo siguiente que no existe. En este último
supuesto, la lista quedaría en estado inconsistente y así el problema se
lo encontraría el próximo hilo que intente usarla.

\subsection{Exclusión mutua}\label{_exclusiuxf3n_mutua}

Para evitar que estas situaciones lleven a la corrupción de datos y a
caídas de servicios y sistemas, debemos asegurarnos que solo un hilo en
cada momento puede manipular recursos y variables compartidas. Por
tanto, necesitamos algún tipo de mecanismo de sincronización para que
mientras se ejecuta \texttt{++count} no se pueda ejecutar
\texttt{-\/-count} en otro hilo, ni viceversa. O para que mientras un
hilo haga un \texttt{insert()} o un \texttt{extract()} en una lista,
otro no pueda utilizar ni estas ni otras funciones de la misma clase.

Para resolver esto, debemos empezar buscando las \textbf{secciones
críticas} de nuestro código. Una \textbf{sección crítica} es una porción
del código donde se accede a variables, tablas, listas, archivos y otros
recursos compartidos.

Para evitar \textbf{condiciones de carrera}, el acceso a las
\textbf{secciones críticas} debe ser controlado, de manera que cuando un
hilo se esté ejecutando en una sección de este tipo ningún otro pueda
hacerlo en la suya correspondiente para manipular los mismos recursos.
En estos casos se dice que existe \textbf{{}exclusión mutua} entre las
\textbf{secciones críticas}.

\subsection{Eventos}\label{_eventos}

Las \textbf{condiciones de carrera} son el principal problema del código
anterior del productor y el consumidor, pero no el único. En ambos
ejemplos se usan bucles para que el hilo espere si el vector está lleno
o vacío, antes de continuar. A esta técnica se la denomina
\textbf{espera ocupada}{} o \textbf{espera activa}{} y está
completamente desaconsejada usarla en código del espacio de usuario,
porque contribuye a gastar tiempo de CPU inútilmente.

En su lugar, se recomienda usar mecanismos de sincronización ofrecidos
por el sistema operativo; diseñados para que un hilo o proceso notifique
eventos a otro, de tal forma que hasta que eso ocurre, el que espera
permanezca en estado \textbf{esperando}, dejando la CPU para los hilos
que la necesitan.

\section{Sincronización por
hardware}\label{_sincronizaciuxf3n_por_hardware}

Las soluciones ofrecidas por el sistema operativo, para resolver los
problemas anteriores, suelen tener que apoyarse en características del
hardware. A continuación veremos algunas de esas características, antes
de profundizar en los mecanismos ofrecidos por el sistema operativo.

\subsection{Bloque de las
interrupciones}\label{_bloque_de_las_interrupciones}

El problema de la sección crítica puede ser resuelto de forma sencilla
en un sistema monoprocesador.

Como el núcleo del sistema operativo es un software controlado mediante
interrupciones, basta con que los hilos bloqueen las interrupciones
mientras se está dentro de la sección crítica. Así, el sistema operativo
no puede tomar el control y asignar otro hilo a la CPU, lo que impide
que se ejecute otra secuencia de instrucciones que podría modificar los
datos compartidos.

Indudablemente esta solución no es práctica en sistema multiprocesador,
donde hay varios procesadores ejecutándose a la vez.

\subsection{Instrucciones atómicas}\label{_instrucciones_atuxf3micas}

{} Las CPU modernas disponen de instrucciones para comparar y modificar
el contenido de una variable o intercambiar el contenido de dos
variables, de forma \textbf{{}atómica}. El término \textbf{atómico} hace
referencia a que las operaciones se ejecutan como una unidad
ininterrumpible. No importa que varias CPU ejecuten estas instrucciones
simultáneamente, puesto el hardware se encargará de que sean ejecutadas
secuencialmente en algún orden arbitrario.

En los procesadores MIPS el acceso atómico a la memoria se realiza
mediante las instrucciones \textbf{\emph{load-linked}} \texttt{ll} y
\textbf{\emph{store-conditional}} \texttt{sc}.

La instrucción \texttt{ll} caga un valor de la memoria en un registro y
monitorizar la dirección de memoria correspondiente para detectar si
otro procesador la modifica. Mientras que la instrucción \texttt{sc}
intenta almacenar un valor en la dirección de memoria, siempre que dicha
dirección no haya sido modificada desde la anterior instrucción
\texttt{ll}.

Ambas instrucciones se utilizan de la siguiente manera para implementar
una operación atómica de incremento:

\begin{verbatim}
main:
  la      r2, count

retry_increment:
  ll      r1, 0(r2)                 
  addiu   r1, r1, 1                 
  sc      r1, 0(r2)                 
  beqz    r1, retry_increment       
  nop                               
\end{verbatim}

\begin{itemize}
\item
  Suponiendo que \texttt{r2} contiene la dirección de la variable
  \texttt{count}, carga el valor de dicha variable en \texttt{r1}.
\item
  Incrementa el valor en el registro \texttt{r1}.
\item
  Intenta almacenar el valor incrementado en \texttt{count}. Si tiene
  éxito, porque el valor de \texttt{count} no ha sido modificado por
  otro proceso, \texttt{r1} se pone a 1. En caso contrario, se pone a 0.
\item
  Si \texttt{r1} vale 0, es que la instrucción \texttt{sc} falló porque
  la variable \texttt{count} fue modificada por otro proceso. En ese
  caso, se salta a \texttt{retry\_increment} para volver a intentar el
  incremento de la variable.
\item
  Se pone una instrucción para evitar que después de \texttt{beqz} se
  ponga un instrucción que el procesador ejecute antes del salto.
\end{itemize}

Antes y después del bloque de instrucciones \texttt{ll} y \texttt{sc},
generalmente se usa la instrucción \texttt{sync}. Esta instrucción crear
una barrera de operaciones en memoria, de forma que el procesador se ve
obligado a completar todas las operaciones en memoria anteriores a
\texttt{sync} antes de realizar las operaciones en memoria siguientes,
como \texttt{ll} y \texttt{sc}.

Estas instrucciones están disponibles para los programadores de C y C++
a través de tipos especiales. Por ejemplo, en C11 \{clang\_stdatomic\}
define tipos como: \texttt{atomic\_bool}, \texttt{atomic\_uint} o
\texttt{atomic\_char} para declarar variables atómicas de los tipos
\texttt{bool}, \texttt{unsigned\ int} y \texttt{char}, respectivamente.
También declara funciones para inicializar, leer, guardar, intercambiar,
sumar, restar y realizar operaciones lógicas, de forma atómica sobre
estas variables:

\begin{Shaded}
\begin{Highlighting}[]
\NormalTok{atomic\_int count}\OperatorTok{;}
\NormalTok{atomic\_init}\OperatorTok{(\&}\NormalTok{count}\OperatorTok{,} \DecValTok{0}\OperatorTok{);} 
\DataTypeTok{int}\NormalTok{ old\_count }\OperatorTok{=}\NormalTok{ atomic\_fetch\_add}\OperatorTok{(\&}\NormalTok{count}\OperatorTok{,} \DecValTok{1}\OperatorTok{);} 
\end{Highlighting}
\end{Shaded}

\begin{itemize}
\item
  Inicializar el valor de la variable atómica.
\item
  Como un \texttt{count++} atómico: incrementa la variable devolviendo
  el valor previo.
\end{itemize}

En C++11, \{cpp\_atomic\_h\} declara la plantilla \{cpp\_atomic\} que
ofrece una funcionada similar:

\begin{Shaded}
\begin{Highlighting}[]
\BuiltInTok{std::}\NormalTok{atomic}\OperatorTok{\textless{}}\DataTypeTok{int}\OperatorTok{\textgreater{}}\NormalTok{ count }\OperatorTok{\{}\DecValTok{0}\OperatorTok{\};}  
\DataTypeTok{int}\NormalTok{ old\_count }\OperatorTok{=}\NormalTok{ count}\OperatorTok{++;}     
\end{Highlighting}
\end{Shaded}

\begin{itemize}
\item
  También hay un tipo \texttt{std::atomic\_int} que es equivalente.
\item
  Se usa el constructor para inicializar la variable atómica.
\item
  Además de soportar los operadores \textquotesingle++\textquotesingle{}
  y \textquotesingle-\/-\textquotesingle, soporta los métodos
  \{cpp\_atomic\_fetch\_add\} y \{cpp\_atomic\_fetch\_sub\} para sumar y
  restar devolviendo el valor previo.
\end{itemize}

La importancia de estas instrucciones está en que pueden ser utilizadas
por el sistema operativo para ofrecer soluciones sencillas al problema
de la sección crítica. Por ejemplo, \textbf{semáforos} o
\textbf{\emph{mutex}}.

\section{Semáforos}\label{_semuxe1foros}

{} La exclusión mutua en las secciones críticas se asegura utilizando
adecuadamente una serie de recursos que para ese fin proporciona el
sistema operativo. Estos recursos utilizan internamente instrucciones y
otras características de la CPU, incluidas por los diseñadores para
resolver este tipo de problemas, que hemos comentado anteriormente. Ese
es el caso de los \textbf{semáforos}.

Los \textbf{semáforos} son un tipo de objetos del sistema operativo que
nos permiten controlar el acceso a una sección crítica, por medio de dos
primitivas: \textbf{acquire} y \textbf{release} ---o \textbf{wait} y
\textbf{signal}, según el libro de texto que consultemos---.

Los \textbf{semáforos} tienen un contador interno que se inicializa, de
tal forma que \textbf{un semáforo inicializado a \emph{N} solo permite
que \emph{N} hilos ejecuten la sección crítica al mismo tiempo}.

A continuación describimos el mecanismo de funcionamiento:

\begin{Shaded}
\begin{Highlighting}[]
\NormalTok{semaphore S}\OperatorTok{(}\DecValTok{10}\OperatorTok{);}    

\NormalTok{S}\OperatorTok{.}\NormalTok{acquire}\OperatorTok{()}         

\CommentTok{// Código de la sección crítica... }

\NormalTok{S}\OperatorTok{.}\NormalTok{signal}\OperatorTok{();}         
\end{Highlighting}
\end{Shaded}

\begin{itemize}
\item
  Crear el \textbf{semáforo} \texttt{S} inicializado a 10. Un
  \textbf{semáforo} contiene fundamentalmente un contador con el número
  máximo de hilos que pueden estar ejecutando el código de la sección
  crítica al mismo tiempo.
\item
  Intentar entrar en la sección crítica:

  \begin{itemize}
  \item
    Si el contador interno del \textbf{semáforo} es mayor que 0,
    \texttt{acquire()} lo decrementa y retorna para que la ejecución
    continúe.
  \item
    Si el contador interno del \textbf{semáforo} es igual a 0,
    \texttt{acquire()} saca al hilo de la CPU y lo pone en una cola de
    espera, suspendiendo así su ejecución. Básicamente, es que hay
    demasiados hilos dentro de la sección crítica.
  \end{itemize}
\item
  Aquí iría el código protegido con el \textbf{semáforo}. Es decir, el
  código de la sección crítica en sí.
\item
  Salir de la sección crítica:

  \begin{itemize}
  \item
    Si el contador interno del \textbf{semáforo} es mayor que 0,
    \texttt{release()} lo incrementa y retorna para que la ejecución
    continúe.
  \item
    Si el contador interno del \textbf{semáforo} es igual a 0,
    \texttt{release()} lo incrementa y saca a uno de los hilos en la
    cola de espera ---donde los puso su \texttt{acquire()}--- para
    meterlo en la cola de preparados, dejándolo listo para entrar en la
    CPU. Cuando eso ocurra, ese hilo decrementará el contador interno
    del \textbf{semáforo} y saldrá de su \texttt{acquire()}, donde hasta
    ahora estaba atrapado. Mientras tanto \texttt{release()} retorna y
    la ejecución del hilo que sale del sección crítica continúa.
  \end{itemize}
\end{itemize}

Para que funcione correctamente, el \textbf{semáforo} \texttt{S} debe
ser el mismo para todos los hilos que tengan secciones críticas en cuya
ejecución debe haber \textbf{exclusión mutua}. Es decir, el
\textbf{semáforo} \texttt{S} debe estar compartido entre los hilos de la
misma manera que las estructuras de datos, variables y otros recursos
que protege.

\subsection{Tipos de semáforos}\label{_tipos_de_semuxe1foros}

Tanto el estándar POSIX como Windows API soportan semáforos y ambos
admiten dos tipos de semáforos:

\begin{itemize}
\item
  Los \textbf{semáforos anónimos}{} que solo existen en el espacio de
  direcciones del proceso que los crea, de tal forma que están
  disponibles para sincronizar hilos del mismo proceso.

  La forma de usarlos para sincronizar procesos diferentes o hilos en
  procesos diferentes depende del sistema operativo. Con Windows API se
  pueden heredar de padres a hijos. Mientras que en sistemas POSIX es
  necesario crear el \textbf{semáforo} en una región de \textbf{memoria
  compartida}, que hayamos creado previamente, e indicar un valor
  distinto de 0 en el argumento \texttt{pshared} de
  \{linux\_sem\_init\}.
\item
  Las \textbf{semáforos con nombre}{} son públicos al resto del sistema,
  por lo que teóricamente cualquier proceso con permisos puede abrirlos
  para utilizarlos.
\end{itemize}

\begin{longtable}[]{@{}
  >{\raggedright\arraybackslash}p{(\linewidth - 4\tabcolsep) * \real{0.3333}}
  >{\raggedright\arraybackslash}p{(\linewidth - 4\tabcolsep) * \real{0.3333}}
  >{\raggedright\arraybackslash}p{(\linewidth - 4\tabcolsep) * \real{0.3333}}@{}}
\caption{Funciones de la API para manipular semáforos.}\tabularnewline
\toprule\noalign{}
\begin{minipage}[b]{\linewidth}\raggedright
\end{minipage} & \begin{minipage}[b]{\linewidth}\raggedright
POSIX API
\end{minipage} & \begin{minipage}[b]{\linewidth}\raggedright
Windows API
\end{minipage} \\
\midrule\noalign{}
\endfirsthead
\toprule\noalign{}
\begin{minipage}[b]{\linewidth}\raggedright
\end{minipage} & \begin{minipage}[b]{\linewidth}\raggedright
POSIX API
\end{minipage} & \begin{minipage}[b]{\linewidth}\raggedright
Windows API
\end{minipage} \\
\midrule\noalign{}
\endhead
\bottomrule\noalign{}
\endlastfoot
\textbf{Crear semáforo anónimo} & \{linux\_sem\_init\} &
\{win32\_createsemaphore\} \\
\textbf{Crear semáforo con nombre} & \{linux\_sem\_open\} &
\{win32\_createsemaphore\} \\
\textbf{Abrir semáforo con nombre} & \{linux\_sem\_open\} &
\{win32\_opensemaphore\} \\
\textbf{Operación acquire} & \{linux\_sem\_wait\} &
\{win32\_waitforsingleobject\} \\
\textbf{Operación release} & \{linux\_sem\_post\} &
\{win32\_releasesemaphore\} \\
\textbf{Cerrar semáforo con nombre} & \{linux\_sem\_close\} &
\{win32\_closehandle\} \\
\textbf{Destruir semáforo anónimo} & \{linux\_sem\_destroy\} &
{{[}Automático{]}} \\
\textbf{Destruir semáforo con nombre} & \{linux\_sem\_unlink\} &
{{[}Automático{]}} \\
\end{longtable}

\subsection{Ejemplos del uso de
semáforos}\label{_ejemplos_del_uso_de_semuxe1foros}

En el ejemplo \{anom\_shared\_memory\_cpp\} de comunicación mediante
memoria compartida, se usa un \textbf{semáforo} para que el proceso hijo
indique al proceso padre que ha terminado de calcular el factorial y el
resultado ya está en la memoria.

En \{shared\_memory\_cpp\} está el ejemplo completo de un programa que
muestra periódicamente la hora del sistema y que puede ser controlado
remotamente, mediante memoria compartida, con un programa como el de
\{shared\_memory\_control\_cpp\}.

Para enviar los mensajes entre ambos programas, en la memoria compartida
se reserva hueco para un búfer, en el que el programa de control copia
el comando que quiere enviar, y para dos \textbf{semáforos}:

\begin{Shaded}
\begin{Highlighting}[]
\KeywordTok{struct}\NormalTok{ memory\_content}
\OperatorTok{\{}
\NormalTok{    sem\_t empty}\OperatorTok{;} 
\NormalTok{    sem\_t ready}\OperatorTok{;} 
    \DataTypeTok{char}\NormalTok{ command\_buffer}\OperatorTok{[}\NormalTok{MAX\_COMMAND\_SIZE}\OperatorTok{];}
\OperatorTok{\};}
\end{Highlighting}
\end{Shaded}

\begin{itemize}
\item
  Indica cuándo \texttt{command\_buffer} está vacío, así que se
  inicializa a 1. El programa de control usa \{linux\_sem\_wait\} en
  este \textbf{semáforo} antes de escribir un nuevo comando en
  \texttt{command\_buffer}:

  \begin{itemize}
  \item
    Si el \textbf{semáforo} está a 0, el programa de control pasa y
    escribe el comando. Después llama a \{linux\_sem\_post\} en
    \texttt{ready}.
  \item
    Si el \textbf{semáforo} está a 1, el programa de control queda
    bloqueado y tiene que esperar a que el programa controlado use
    \{linux\_sem\_post\} sobre el mismo \textbf{semáforo}. El programa
    controlado lo hace después de leer el comando para interpretarlo.
  \end{itemize}
\item
  Indica cuándo \texttt{command\_buffer} tiene un comando, así que se
  inicializa a 0. El programa controlado usa \{linux\_sem\_wait\} en
  este \textbf{semaforo} antes de leer el comando en
  \texttt{command\_buffer} para interpretarlo:

  \begin{itemize}
  \item
    Si el \textbf{semáforo} está a 0, el programa de control pasa y lee
    el comando. Después llama a \{linux\_sem\_post\} en \texttt{empty}.
  \item
    Si el \textbf{semáforo} está a 1, el programa controlado queda
    bloqueado y tiene que esperar a que el programa de control use
    \{linux\_sem\_post\} sobre el mismo \textbf{semáforo}. El programa
    de control lo hace después de escribir un nuevo comando en
    \texttt{command\_buffer}.
  \end{itemize}
\end{itemize}

El detalle de cómo ambos programas usan los semáforos, se puede ver en
el código de \{shared\_memory\_control\_cpp\} y \{shared\_memory\_cpp\},
respectivamente.

Finalmente, para resolver el \textbf{problema del productor-consumidor}
tenemos que considerar que:

\begin{itemize}
\item
  Necesitamos un semáforo para que haya \textbf{exclusión mutua} entre
  ambos al insertar y extraer elementos del vector.
\item
  Necesitamos una forma de que el productor espere cuando el vector esté
  lleno y que el consumidor haga lo mismo cuando el vector esté vacío.
  Una solución es usar dos semáforos, uno para que cuente el número de
  elementos en el vector y otro para contar el número de huecos libres:
\end{itemize}

\begin{Shaded}
\begin{Highlighting}[]
\DataTypeTok{sem\_t}\NormalTok{ mutex}\OperatorTok{;}       
\DataTypeTok{sem\_t}\NormalTok{ fill\_count}\OperatorTok{;}  
\DataTypeTok{sem\_t}\NormalTok{ empty\_count}\OperatorTok{;} 

\BuiltInTok{std::}\NormalTok{vector}\OperatorTok{\textless{}}\DataTypeTok{item\_t}\OperatorTok{\textgreater{}}\NormalTok{ buffer}\OperatorTok{(}\NormalTok{ VECTOR\_SIZE }\OperatorTok{);}

\DataTypeTok{void}\NormalTok{ initialize}\OperatorTok{()}
\OperatorTok{\{}
\NormalTok{    sem\_init}\OperatorTok{(} \OperatorTok{\&}\NormalTok{mutex}\OperatorTok{,} \DecValTok{0}\OperatorTok{,} \DecValTok{1} \OperatorTok{);}  
\NormalTok{    sem\_init}\OperatorTok{(} \OperatorTok{\&}\NormalTok{mutex}\OperatorTok{,} \DecValTok{0}\OperatorTok{,} \DecValTok{0} \OperatorTok{);}  
\NormalTok{    sem\_init}\OperatorTok{(} \OperatorTok{\&}\NormalTok{mutex}\OperatorTok{,} \DecValTok{0}\OperatorTok{,}\NormalTok{ VECTOR\_SIZE }\OperatorTok{);}  
\OperatorTok{\}}

\DataTypeTok{void}\NormalTok{ productor}\OperatorTok{()}
\OperatorTok{\{}
    \CommentTok{// ...}

    \ControlFlowTok{while}\OperatorTok{(}\CommentTok{/* ... */}\OperatorTok{)}
    \OperatorTok{\{}
        \DataTypeTok{item\_t}\NormalTok{ item }\OperatorTok{=}\NormalTok{ produce\_item}\OperatorTok{();}

\NormalTok{        sem\_wait}\OperatorTok{(} \OperatorTok{\&}\NormalTok{empty\_count }\OperatorTok{);} 
\NormalTok{        sem\_wait}\OperatorTok{(} \OperatorTok{\&}\NormalTok{mutex }\OperatorTok{);}       

\NormalTok{        vector}\OperatorTok{.}\NormalTok{push\_back}\OperatorTok{(}\NormalTok{item}\OperatorTok{);}

\NormalTok{        sem\_post}\OperatorTok{(} \OperatorTok{\&}\NormalTok{mutex }\OperatorTok{);}       
\NormalTok{        sem\_post}\OperatorTok{(} \OperatorTok{\&}\NormalTok{fill\_count }\OperatorTok{);}  
    \OperatorTok{\}}

    \CommentTok{// ...}
\OperatorTok{\}}

\DataTypeTok{void}\NormalTok{ consumer}\OperatorTok{()}
\OperatorTok{\{}
    \CommentTok{// ...}

    \ControlFlowTok{while}\OperatorTok{(}\CommentTok{/* ... */}\OperatorTok{)}
    \OperatorTok{\{}
\NormalTok{        sem\_wait}\OperatorTok{(} \OperatorTok{\&}\NormalTok{fill\_count }\OperatorTok{);}  
\NormalTok{        sem\_wait}\OperatorTok{(} \OperatorTok{\&}\NormalTok{mutex }\OperatorTok{);}       

        \DataTypeTok{item\_t}\NormalTok{ item }\OperatorTok{=}\NormalTok{ vector}\OperatorTok{.}\NormalTok{back}\OperatorTok{();}
\NormalTok{        vector}\OperatorTok{.}\NormalTok{pop\_back}\OperatorTok{();}

\NormalTok{        sem\_post}\OperatorTok{(} \OperatorTok{\&}\NormalTok{mutex }\OperatorTok{);}       
\NormalTok{        sem\_post}\OperatorTok{(} \OperatorTok{\&}\NormalTok{empty\_count }\OperatorTok{);} 

\NormalTok{        consume\_item}\OperatorTok{(}\NormalTok{item}\OperatorTok{);}
    \OperatorTok{\}}

    \CommentTok{// ...}
\OperatorTok{\}}
\end{Highlighting}
\end{Shaded}

\begin{itemize}
\item
  \textbf{Semáforo} que se encarga de la exclusión mutua. Se inicializa
  a 1, para que el primer hilo que use \{linux\_sem\_wait\} pueda
  entrar.
\item
  \textbf{Semáforo} que se encarga de contar huecos ocupados en el
  vector. Se inicializa a 0, porque al principio no hay ningún elemento.
\item
  \textbf{Semáforo} que se encarga de contar los huecos libres en el
  vector. Se inicializa a VECTOR\_SIZE, porque están todos vacíos.
\item
  El segundo argumento de \{linux\_sem\_init\} es \texttt{pshared}. Se
  pone a 0 para indicar que este \textbf{semaforo} no se va a compartir
  entre procesos diferentes.
\item
  Antes de insertar un elemento se decrementa \texttt{empty\_count}.
  Así, si vale 0, es que el vector está lleno y el productor se bloquea.
  El consumidor incrementa \texttt{empty\_count} tras extraer un
  elemento y dejar hueco, despertando al productor.
\item
  Antes de extraer un elemento se decrementa \texttt{fill\_count}. Así,
  si vale 0, es que el vector está vacío y el consumidor se bloquea. El
  productor incrementa \texttt{fill\_count} tras insertar un elemento
  nuevo, despertando al consumidor.
\item
  El acceso al vector con los elementos es en \textbf{exclusión mutua},
  así que tanto productor como consumidor deben usar
  \{linux\_sem\_wait\} sobre \texttt{mutex} antes de acceder a él. Esto
  decrementa el semáforo, así que solo uno de los dos pasa y ejecuta las
  líneas siguientes, mientras el otro queda bloqueado. Cuando el que
  haya pasado termine, debe usar \{linux\_sem\_post\} sobre
  \texttt{mutex} para incrementar el semáforo y permitir que el otro
  hilo entre en su \textbf{sección crítica}.
\end{itemize}

\section{Mutex}\label{_mutex}

Los \textbf{\emph{{}mutex}} ---término que tiene su origen en
\emph{mutual exclusion}--- son un tipo de objeto del sistema operativo
que permite controlar el acceso a una sección crítica, de forma que
\textbf{solo un hilo pueda ejecutarla al mismo tiempo}.

En este sentido, los \textbf{\emph{mutex}} se comportan como semáforos
inicializados a 1, motivo por el que también se denominan
\textbf{semáforos binarios}{}. Por tanto, aunque un sistema o lenguaje
solo soporte \textbf{semáforos}, es directo usarlos para implementar
\textbf{\emph{mutex}}.

El estándar POSIX soporta \textbf{\emph{mutex}} a través de
\{pthreads\}. Por defecto solo se pueden utilizar para sincronizar hilos
del mismo proceso; pero tienen un atributo para permitir la
sincronización entre procesos diferentes, aunque para eso deben ser
creados en una región de memoria compartida por dichos procesos.

La API de \{pthreads\} no soporta \textbf{semáforos} porque, como vimos
antes, ya eran parte del estándar POSIX.

\begin{longtable}[]{@{}
  >{\raggedright\arraybackslash}p{(\linewidth - 8\tabcolsep) * \real{0.2000}}
  >{\raggedright\arraybackslash}p{(\linewidth - 8\tabcolsep) * \real{0.2000}}
  >{\raggedright\arraybackslash}p{(\linewidth - 8\tabcolsep) * \real{0.2000}}
  >{\raggedright\arraybackslash}p{(\linewidth - 8\tabcolsep) * \real{0.2000}}
  >{\raggedright\arraybackslash}p{(\linewidth - 8\tabcolsep) * \real{0.2000}}@{}}
\caption{Funciones de la API para manipular
\emph{mutex}.}\tabularnewline
\toprule\noalign{}
\begin{minipage}[b]{\linewidth}\raggedright
\end{minipage} & \begin{minipage}[b]{\linewidth}\raggedright
C++
\end{minipage} & \begin{minipage}[b]{\linewidth}\raggedright
POSIX API
\end{minipage} &
\multicolumn{2}{>{\raggedright\arraybackslash}p{(\linewidth - 8\tabcolsep) * \real{0.4000} + 2\tabcolsep}@{}}{\begin{minipage}[b]{\linewidth}\raggedright
Windows API
\end{minipage} \\
\midrule\noalign{}
\endfirsthead
\toprule\noalign{}
\begin{minipage}[b]{\linewidth}\raggedright
\end{minipage} & \begin{minipage}[b]{\linewidth}\raggedright
C++
\end{minipage} & \begin{minipage}[b]{\linewidth}\raggedright
POSIX API
\end{minipage} &
\multicolumn{2}{>{\raggedright\arraybackslash}p{(\linewidth - 8\tabcolsep) * \real{0.4000} + 2\tabcolsep}@{}}{\begin{minipage}[b]{\linewidth}\raggedright
Windows API
\end{minipage} \\
\midrule\noalign{}
\endhead
\bottomrule\noalign{}
\endlastfoot
\textbf{Crear} & \{cpp\_mutex\} & \{pthread\_mutex\_init\} &
\{win32\_initializecriticalsection\} & \{win32\_createmutex\} \\
\textbf{Abrir} & & & & \{win32\_openmutex\} \\
\textbf{Operación acquire} & \{cpp\_mutex\_lock\} &
\{pthread\_mutex\_lock\} & \{win32\_entercriticalsection\} &
\{win32\_waitforsingleobject\} \\
\textbf{Operación release} & \{cpp\_mutex\_unlock\} &
\{pthread\_mutex\_unlock\} & \{win32\_leavecriticalsection\} &
\{win32\_releasemutex\} \\
\textbf{Destruir} & {{[}Destructor{]}} & \{pthread\_mutex\_destroy\} &
\{win32\_deletecriticalsection\} & \{win32\_closehandle\} \\
\end{longtable}

En Windows API hay dos tipo de objetos equiparables a los
\textbf{\emph{mutex}}: los \textbf{\emph{mutex}} y las \textbf{secciones
críticas}. Las \textbf{secciones críticas} son más ligeras, pero solo se
pueden utilizar para sincronizar hilos del mismo proceso. Mientras que
los \textbf{\emph{mutex}} de Windows API son objetos más costosos, pero
se pueden compartir entre procesos sin utilizar memoria compartida; ya
sea mediante herencia al crear un proceso hijo o asignando un nombre al
\textbf{\emph{mutex}}, como ocurre con los \textbf{semáforos}.

\subsection{Ejemplos del uso de mutex}\label{_ejemplos_del_uso_de_mutex}

En \{pthreads\_sync\_factorial\_cpp\} se puede estudiar el código
completo de un ejemplo similar a \{pthreads\_cpp\}, donde se calculaba
el factorial de un número, repartiendo la tarea entre dos hilos, usando
la API de \{pthreads\}. La diferencia es que ahora los hilos no retornan
el resultado, sino que cada uno lo mete en un vector compartido. Al
terminar, el hilo principal recorre el vector multiplicando los
resultados parciales.

Como ahora ambos hilos acceden a una estructura de datos compartida,
esta debe ir protegida por un \textbf{\emph{mutex}}:

\begin{Shaded}
\begin{Highlighting}[]
\DataTypeTok{pthread\_mutex\_t}\NormalTok{ mutex}\OperatorTok{;}
\BuiltInTok{std::}\NormalTok{vector}\OperatorTok{\textless{}}\NormalTok{BigInt}\OperatorTok{\textgreater{}}\NormalTok{ partials}\OperatorTok{;}
\end{Highlighting}
\end{Shaded}

Antes de meter un nuevo valor, cada hilo debe adquirir el
\texttt{mutex}:

\begin{Shaded}
\begin{Highlighting}[]
\CommentTok{// Bloquear el mutex y guardar el resultado}
\NormalTok{pthread\_mutex\_lock}\OperatorTok{(}\NormalTok{ mutex }\OperatorTok{);}   
\NormalTok{partials}\OperatorTok{.}\NormalTok{push\_back}\OperatorTok{(}\NormalTok{ result }\OperatorTok{);}
\NormalTok{pthread\_mutex\_unlock}\OperatorTok{(}\NormalTok{ mutex }\OperatorTok{);} 
\end{Highlighting}
\end{Shaded}

\begin{itemize}
\item
  Adquirir \texttt{mutex} antes de entrar en la \textbf{sección
  crítica}.
\item
  Liberar \texttt{mutex} para salir de la \textbf{sección crítica}. Es
  importante no olvidarnos de liberar el \textbf{\emph{mutex}} al
  terminar o de lo contrario uno de los hilos quedará dormido
  indefinidamente, al no poder entrar en la \textbf{sección crítica}.
\end{itemize}

En \{pthreads\_sync\_counter\_cpp\} se puede observar un ejemplo más
simple donde dos hilos intentan incrementar un contador. El resultado es
correcto siempre que cada hilo adquiera el mismo \textbf{\emph{mutex}}
antes del incremento:

\begin{Shaded}
\begin{Highlighting}[]
\CommentTok{// Bloquear el mutex antes de incrementar el contador.}
\NormalTok{pthread\_mutex\_lock}\OperatorTok{(} \OperatorTok{\&}\NormalTok{args}\OperatorTok{{-}\textgreater{}}\NormalTok{mutex }\OperatorTok{);}
\NormalTok{args}\OperatorTok{{-}\textgreater{}}\NormalTok{counter}\OperatorTok{++;}
\CommentTok{// Desbloquear el mutex tras incrementar el contador.}
\NormalTok{pthread\_mutex\_unlock}\OperatorTok{(} \OperatorTok{\&}\NormalTok{args}\OperatorTok{{-}\textgreater{}}\NormalTok{mutex }\OperatorTok{);}
\end{Highlighting}
\end{Shaded}

En \{threads\_sync\_counter\_cpp\} y \{threads\_sync\_factorial\_cpp\}
se puede ver ejemplos equivalentes pero usando \{cpp\_thread\} y
\{cpp\_mutex\}, de la librería estándar de C++.

\section{Variables de condición}\label{_variables_de_condiciuxf3n}

{} En la solución que dimos al \textbf{problema del
productor-consumidor} usando \textbf{semáforos} (véase el
\hyperref[_ejemplos_del_uso_de_semuxe1foros]{Ejemplos del uso de
semáforos}) empleamos \textbf{semáforos} para implementar las esperas
del productor y el consumidor cuando el vector está lleno o vacío,
respectivamente. Lamentablemente, los \emph{mutex} no se pueden usar de
la misma manera para señalar eventos. En su lugar necesitamos otro tipo
de objeto llamado \textbf{variable de condición}.

Las \textbf{variables de condición} soportan tres primitivas
principales:

\begin{description}
\item[wait( mutex )]
Es llamada por un hilo que desea esperar a que ocurra el evento que
representa la variable de condición. El hilo debe haber adquirido antes
el \textbf{\emph{mutex}}, es liberado en el momento de poner al hilo en
estado \textbf{esperando}. Varios hilos pueden llamar a \textbf{wait}
sobre la misma variable de condición, a la espera de que alguno use
\textbf{notify}.
\item[notify]
Es llamada por un hilo que quiere notificar el suceso de un evento a los
hilos que esperan en la variable de condición. Uno de esos hilos es
despertado, adquiere el \textbf{\emph{mutex}} que liberó al llamar a
\textbf{wait} y, finalmente, retorna de \textbf{wait} para seguir
ejecutándose.
\item[notifyAll]
Es llamada por un hilo que quiere notificar el suceso de un evento a los
hilos que esperan en la variable de condición. Todos los hilos son
despertados e intentan adquirir el \textbf{\emph{mutex}} que liberaron
al llamar a \textbf{wait}. Cuando lo consiguen, retornan de
\textbf{wait} para seguir ejecutándose. Obviamente, si todos hicieron
\textbf{wait} sobre el mismo \textbf{\emph{mutex}}, irán retornando de
\textbf{wait} de uno en uno, porque solo un hilo puede tener el
\textbf{\emph{mutex}} al mismo tiempo.
\end{description}

Tanto Windows API como el estándar POSIX, a través de \{pthreads\},
soportan \textbf{variables de condición}.

\begin{longtable}[]{@{}
  >{\raggedright\arraybackslash}p{(\linewidth - 6\tabcolsep) * \real{0.2500}}
  >{\raggedright\arraybackslash}p{(\linewidth - 6\tabcolsep) * \real{0.2500}}
  >{\raggedright\arraybackslash}p{(\linewidth - 6\tabcolsep) * \real{0.2500}}
  >{\raggedright\arraybackslash}p{(\linewidth - 6\tabcolsep) * \real{0.2500}}@{}}
\caption{Funciones de la API para manipular variables de
condición.}\tabularnewline
\toprule\noalign{}
\begin{minipage}[b]{\linewidth}\raggedright
\end{minipage} & \begin{minipage}[b]{\linewidth}\raggedright
C++
\end{minipage} & \begin{minipage}[b]{\linewidth}\raggedright
POSIX API
\end{minipage} & \begin{minipage}[b]{\linewidth}\raggedright
Windows API
\end{minipage} \\
\midrule\noalign{}
\endfirsthead
\toprule\noalign{}
\begin{minipage}[b]{\linewidth}\raggedright
\end{minipage} & \begin{minipage}[b]{\linewidth}\raggedright
C++
\end{minipage} & \begin{minipage}[b]{\linewidth}\raggedright
POSIX API
\end{minipage} & \begin{minipage}[b]{\linewidth}\raggedright
Windows API
\end{minipage} \\
\midrule\noalign{}
\endhead
\bottomrule\noalign{}
\endlastfoot
\textbf{Crear} & \{cpp\_condition\_variable\} & \{pthread\_cond\_init\}
& \{win32\_initializeconditionvariable\} \\
\textbf{Operación wait} & \{cpp\_condition\_variable\_wait\} &
\{pthread\_cond\_wait\} & \{win32\_sleepconditionvariablecs\} \\
\textbf{Operación notify} & \{cpp\_condition\_variable\_notify\_one\} &
\{pthread\_cond\_signal\} & \{win32\_wakeconditionvariable\} \\
\textbf{Operación notifyAll} & \{cpp\_condition\_variable\_notify\_all\}
& \{pthread\_cond\_broadcast\} & \{win32\_wakeallconditionvariable\} \\
\textbf{Destruir} & {{[}Destructor{]}} & \{pthread\_cond\_destroy\} & \\
\end{longtable}

Por defecto, las \textbf{variables de condición} de \{pthreads\} solo se
pueden utilizar para sincronizar hilos del mismo proceso; pero tienen un
atributo para permitir la sincronización entre procesos diferentes.
Obviamente, para eso deben ser creadas en una región de memoria
compartida por dichos procesos.

Las \textbf{variables de condición} de Windows API solo se pueden
utilizar en hilos del mismo procesos. Como alternativa, Windows API
soporta \textbf{eventos}, que son un tipo de objeto similar a las
\textbf{variables de condición}, pero que sí se puede utilizar entre
hilos de procesos diferentes\hyperref[Win32-EventObjects]{???}.

A los \textbf{eventos} se les puede asignar un nombre, para que sean
accesibles por otros procesos, o heredarse de padres a hijos. Además son
más pesados que las \textbf{variables de condición} de Windows API, no
exigen un \textbf{\emph{mutex}} para liberar al invocar su operación
\textbf{wait}, ni admiten la operación \textbf{notifyAll}.

\begin{longtable}[]{@{}
  >{\raggedright\arraybackslash}p{(\linewidth - 2\tabcolsep) * \real{0.5000}}
  >{\raggedright\arraybackslash}p{(\linewidth - 2\tabcolsep) * \real{0.5000}}@{}}
\caption{Funciones de la API para manipular eventos de Windows
API.}\tabularnewline
\toprule\noalign{}
\begin{minipage}[b]{\linewidth}\raggedright
\end{minipage} & \begin{minipage}[b]{\linewidth}\raggedright
Windows API
\end{minipage} \\
\midrule\noalign{}
\endfirsthead
\toprule\noalign{}
\begin{minipage}[b]{\linewidth}\raggedright
\end{minipage} & \begin{minipage}[b]{\linewidth}\raggedright
Windows API
\end{minipage} \\
\midrule\noalign{}
\endhead
\bottomrule\noalign{}
\endlastfoot
\textbf{Crear} & \{win32\_createevent\} \\
\textbf{Abrir} & \{win32\_openevent\} \\
\textbf{Operación wait} & \{win32\_waitforsingleobject\} \\
\textbf{Operación notify} & \{win32\_setevent\} \\
\textbf{Resetear evento} & \{win32\_resetevent\} \\
\textbf{Destruir} & \{win32\_closehandle\} \\
\end{longtable}

\subsection{Ejemplos del uso de variables de
condición}\label{_ejemplos_del_uso_de_variables_de_condiciuxf3n}

Vamos a intentar resolver el \textbf{problema del productor-consumidor}
sin usar \textbf{semáforos}. Para lo que, nuevamente, tenemos que
considerar que:

\begin{itemize}
\item
  Necesitamos \textbf{exclusión mutua} entre ambos hilos al insertar y
  extraer elementos del vector para evitar \textbf{condiciones de
  carrera}, por tanto usamos un \textbf{\emph{mutex}} para proteger la
  \textbf{sección crítica}.
\item
  Necesitamos una forma de que el productor espere cuando el vector está
  lleno y que el consumidor haga lo mismo cuando el vector está vacío.
  Para señalar estos eventos necesitamos dos \textbf{variables de
  condición}:
\end{itemize}

\begin{Shaded}
\begin{Highlighting}[]
\BuiltInTok{std::}\NormalTok{mutext mutex}\OperatorTok{;}                
\BuiltInTok{std::}\NormalTok{condition\_variable no\_full}\OperatorTok{;}  
\BuiltInTok{std::}\NormalTok{condition\_variable no\_empty}\OperatorTok{;} 

\BuiltInTok{std::}\NormalTok{vector}\OperatorTok{\textless{}}\DataTypeTok{item\_t}\OperatorTok{\textgreater{}}\NormalTok{ buffer}\OperatorTok{(}\NormalTok{ VECTOR\_SIZE }\OperatorTok{);}

\DataTypeTok{void}\NormalTok{ productor}\OperatorTok{()}
\OperatorTok{\{}
    \CommentTok{// ...}

    \ControlFlowTok{while}\OperatorTok{(}\CommentTok{/* ... */}\OperatorTok{)}
    \OperatorTok{\{}
        \DataTypeTok{item\_t}\NormalTok{ item }\OperatorTok{=}\NormalTok{ produce\_item}\OperatorTok{();}

        \BuiltInTok{std::}\NormalTok{unique\_lock lock}\OperatorTok{\{}\NormalTok{ mutex }\OperatorTok{\};}        

        \ControlFlowTok{while}\OperatorTok{(}\NormalTok{ buffer}\OperatorTok{.}\NormalTok{size}\OperatorTok{()} \OperatorTok{==}\NormalTok{ VECTOR\_SIZE }\OperatorTok{)}  
        \OperatorTok{\{}
\NormalTok{            no\_full}\OperatorTok{.}\NormalTok{wait}\OperatorTok{(}\NormalTok{ mutex }\OperatorTok{);}            
        \OperatorTok{\}}

\NormalTok{        vector}\OperatorTok{.}\NormalTok{push\_back}\OperatorTok{(}\NormalTok{ item }\OperatorTok{);}             

\NormalTok{        no\_empty}\OperatorTok{.}\NormalTok{notify\_one}\OperatorTok{();}                
    \OperatorTok{\}}                                         

    \CommentTok{// ...}
\OperatorTok{\}}

\DataTypeTok{void}\NormalTok{ consumer}\OperatorTok{()}
\OperatorTok{\{}
    \CommentTok{// ...}

    \ControlFlowTok{while}\OperatorTok{(}\CommentTok{/* ... */}\OperatorTok{)}
    \OperatorTok{\{}
        \BuiltInTok{std::}\NormalTok{unique\_lock lock}\OperatorTok{\{}\NormalTok{ mutex }\OperatorTok{\};}  

        \ControlFlowTok{while}\OperatorTok{(}\NormalTok{ buffer}\OperatorTok{.}\NormalTok{size}\OperatorTok{()} \OperatorTok{==} \DecValTok{0} \OperatorTok{)}      
        \OperatorTok{\{}
\NormalTok{            no\_empty}\OperatorTok{.}\NormalTok{wait}\OperatorTok{(}\NormalTok{ mutex }\OperatorTok{);}     
        \OperatorTok{\}}

        \DataTypeTok{item\_t}\NormalTok{ item }\OperatorTok{=}\NormalTok{ vector}\OperatorTok{.}\NormalTok{back}\OperatorTok{();}    
\NormalTok{        vector}\OperatorTok{.}\NormalTok{pop\_back}\OperatorTok{();}

\NormalTok{        no\_full}\OperatorTok{.}\NormalTok{notify\_one}\OperatorTok{();}           

\NormalTok{        consume\_item}\OperatorTok{(}\NormalTok{item}\OperatorTok{);}
    \OperatorTok{\}}                                   

    \CommentTok{// ...}
\OperatorTok{\}}
\end{Highlighting}
\end{Shaded}

\begin{itemize}
\item
  \textbf{\emph{Mutex}} que se encarga de la exclusión mutua.
\item
  \textbf{Variable de condición} que se encarga de indicar cuando el
  vector no está lleno.
\item
  \textbf{Variable de condición} que se encarga de indicar cuando el
  vector no está vacío.
\item
  Antes de acceder al vector es necesario bloquear \texttt{mutex}. Ni
  siquiera es seguro preguntar por el número de elementos guardados en
  \texttt{vector} sin antes adquirir el \textbf{\emph{mutex}}, puesto
  que el otro hilo puede estar modificando \texttt{vector} al mismo
  tiempo.
\item
  Los \textbf{\emph{mutex}} se pueden adquirir y liberar con
  \{cpp\_mutex\_lock\} y cpp\_mutex\_unlock\}, pero esa no es la forma
  recomendada. Lo recomendado es crear alguno de los objetos \emph{lock}
  incluidos en la librería estándar. Estos objetos bloquean el
  \textbf{\emph{mutex}} al crearse y lo desbloquean automáticamente al
  destruirse. Así es complicado que nos olvidemos de desbloquearlo al
  salir de la función.
\item
  Antes de insertar un elemento se comprueba si hay algún hueco
  disponible. Si no lo hay, se pone el hilo a la espera en la
  \textbf{variable de condición} \texttt{no\_full}.
\item
  El consumidor despierta al productor de esa espera tras extraer un
  elemento, porque es seguro que al hacerlo habrá dejado un hueco.
\item
  Antes de extraer un elemento se comprueba si hay alguno en el vector.
  Si no lo hay, se pone el hilo a la espera en la \textbf{variable de
  condición} \texttt{no\_empty}.
\item
  El productor despierta al consumidor de esta espera tras insertar un
  nuevo elemento.
\item
  El estándar de C++ indica que las esperas en las \textbf{variables de
  condición} son susceptibles de despertar de forma espuria. Es decir,
  que el hilo puede salir de \{cpp\_condition\_variable\_wait\} sin que
  haya habido notificación. Por eso hay que volver a comprobar la
  condición antes de continuar ejecutando sentencias en la
  \textbf{sección crítica}.
\item
  Cuando los hilos se bloquean en \{cpp\_condition\_variable\_wait\},
  \texttt{mutex} es liberado para que el otro hilo pueda entrar y
  extraer o insertar un elemento. De lo contrario, no podría hacerlo y
  ambos se quedarían bloqueados indefinidamente ---en una situación que
  se denomina \textbf{{}interbloqueo} o \textbf{\emph{deadlock}}---.
  Pero antes de salir de \{cpp\_condition\_variable\_wait\}, el hilo
  adquiere de nuevo el \texttt{mutex}. Así que el código que inserta y
  extrae elementos se ejecuta en \textbf{exclusión mutua}, tal y como
  nos interesa.
\end{itemize}

\section{Implementación de
semáforos}\label{_implementaciuxf3n_de_semuxe1foros}

Como hemos visto, tanto los \textbf{\emph{mutex}} como las
\textbf{variables de condición} son casos particulares de los
\textbf{semáforos}. De igual forma, si un sistema o lenguaje soporta
\textbf{\emph{mutex}} y \textbf{variables de condición}, es muy sencillo
implementar \textbf{semáforos}, así como otras primitivas de
sincronización.

Un \textbf{semáforo} se implementa de forma muy parecida a la solución
anterior al \textbf{problema del productor-consumidor} (véase el
\hyperref[_ejemplos_del_uso_de_variables_de_condiciuxf3n]{Ejemplos del
uso de variables de condición}), solo que eliminando lo relacionado con
el vector de elementos.

El código fuente completo de la clase \texttt{semaphore} está disponible
en \{semaphore\_hpp\}.

\begin{Shaded}
\begin{Highlighting}[]
\KeywordTok{class}\NormalTok{ semaphore }\OperatorTok{\{}
\KeywordTok{public}\OperatorTok{:}

\NormalTok{    semaphore}\OperatorTok{(}\DataTypeTok{int}\NormalTok{ count }\OperatorTok{=} \DecValTok{0}\OperatorTok{)} \OperatorTok{:} \VariableTok{count\_}\OperatorTok{(}\NormalTok{count}\OperatorTok{)} \OperatorTok{\{\}}

    \DataTypeTok{void}\NormalTok{ release}\OperatorTok{()} \OperatorTok{\{}
        \BuiltInTok{std::}\NormalTok{unique\_lock}\OperatorTok{\textless{}}\BuiltInTok{std::}\NormalTok{mutex}\OperatorTok{\textgreater{}}\NormalTok{ lock}\OperatorTok{(}\VariableTok{mutex\_}\OperatorTok{);} 
        \VariableTok{count\_}\OperatorTok{++;}
        \VariableTok{cv\_}\OperatorTok{.}\NormalTok{notify\_one}\OperatorTok{();}                          
    \OperatorTok{\}}

    \DataTypeTok{void}\NormalTok{ acquire}\OperatorTok{()} \OperatorTok{\{}
        \BuiltInTok{std::}\NormalTok{unique\_lock}\OperatorTok{\textless{}}\BuiltInTok{std::}\NormalTok{mutex}\OperatorTok{\textgreater{}}\NormalTok{ lock}\OperatorTok{(}\VariableTok{mutex\_}\OperatorTok{);} 
        \ControlFlowTok{while}\OperatorTok{(}\VariableTok{count\_} \OperatorTok{==} \DecValTok{0}\OperatorTok{)} \OperatorTok{\{}
            \VariableTok{cv\_}\OperatorTok{.}\NormalTok{wait}\OperatorTok{(}\NormalTok{lock}\OperatorTok{);}                        
        \OperatorTok{\}}
        \VariableTok{count\_}\OperatorTok{{-}{-};}
    \OperatorTok{\}}

\KeywordTok{private}\OperatorTok{:}
    \BuiltInTok{std::}\NormalTok{mutex }\VariableTok{mutex\_}\OperatorTok{;}
    \BuiltInTok{std::}\NormalTok{condition\_variable }\VariableTok{cv\_}\OperatorTok{;}
    \DataTypeTok{int} \VariableTok{count\_}\OperatorTok{;}
\OperatorTok{\};}
\end{Highlighting}
\end{Shaded}

\begin{itemize}
\item
  El acceso al contador \texttt{count\_} esta protegido por el
  \textbf{\emph{mutex}} \texttt{mutex\_}.
\item
  En \texttt{semaphore::acquire()}, si el contador está a 0 no pueden
  entrar más hilos, por lo que el hilo debe esperar en la
  \textbf{variable de condición} \texttt{cv\_}.
\item
  Cada vez que un hilo llama a \texttt{semaphore::release()} se
  incrementa el contador y se despierta uno de los hilos a la espera en
  \texttt{cv\_}.
\end{itemize}

En \{threads\_sync\_semaphore\_cpp\} se puede ver un ejemplo de cómo
utilizarla para limitar el número de hilos que pueden ejecutarse en una
\textbf{sección crítica} al mismo tiempo.

\begin{Shaded}
\begin{Highlighting}[]
\DataTypeTok{void}\NormalTok{ thread\_function}\OperatorTok{(}\NormalTok{semaphore}\OperatorTok{\&}\NormalTok{ sem}\OperatorTok{,} \DataTypeTok{int}\NormalTok{ thread\_id}\OperatorTok{)}
\OperatorTok{\{}
\NormalTok{    sem}\OperatorTok{.}\NormalTok{acquire}\OperatorTok{();} 
    \BuiltInTok{std::}\NormalTok{println}\OperatorTok{(} \StringTok{"Hilo }\SpecialCharTok{\{\}}\StringTok{ creado"}\OperatorTok{,}\NormalTok{ thread\_id }\OperatorTok{);}

    \CommentTok{// Dormir el hilo para simular trabajo}
    \CommentTok{// ...}

    \BuiltInTok{std::}\NormalTok{println}\OperatorTok{(} \StringTok{"Hilo }\SpecialCharTok{\{\}}\StringTok{ terminado"}\OperatorTok{,}\NormalTok{ thread\_id }\OperatorTok{);}
\NormalTok{    sem}\OperatorTok{.}\NormalTok{release}\OperatorTok{();} 
\OperatorTok{\}}

\DataTypeTok{int}\NormalTok{ main}\OperatorTok{()}
\OperatorTok{\{}
    \AttributeTok{const} \DataTypeTok{int}\NormalTok{ num\_threads }\OperatorTok{=} \DecValTok{10}\OperatorTok{;}
\NormalTok{    semaphore sem}\OperatorTok{(}\DecValTok{3}\OperatorTok{);} 

    \BuiltInTok{std::}\NormalTok{vector}\OperatorTok{\textless{}}\BuiltInTok{std::}\NormalTok{thread}\OperatorTok{\textgreater{}}\NormalTok{ threads}\OperatorTok{;}
    \ControlFlowTok{for}\OperatorTok{(}\DataTypeTok{int}\NormalTok{ i }\OperatorTok{=} \DecValTok{0}\OperatorTok{;}\NormalTok{ i }\OperatorTok{\textless{}}\NormalTok{ num\_threads}\OperatorTok{;}\NormalTok{ i}\OperatorTok{++)} 
    \OperatorTok{\{}
\NormalTok{        threads}\OperatorTok{.}\NormalTok{push\_back}\OperatorTok{(} \BuiltInTok{std::}\NormalTok{thread}\OperatorTok{(}\NormalTok{ thread\_function}\OperatorTok{,} \BuiltInTok{std::}\NormalTok{ref}\OperatorTok{(}\NormalTok{sem}\OperatorTok{),}\NormalTok{ i }\OperatorTok{));}
    \OperatorTok{\}}

    \ControlFlowTok{for}\OperatorTok{(}\KeywordTok{auto}\OperatorTok{\&}\NormalTok{ t }\OperatorTok{:}\NormalTok{ threads}\OperatorTok{)}
    \OperatorTok{\{}
\NormalTok{        t}\OperatorTok{.}\NormalTok{join}\OperatorTok{();}
    \OperatorTok{\}}

    \ControlFlowTok{return}\NormalTok{ EXIT\_SUCCESS}\OperatorTok{;}
\OperatorTok{\}}
\end{Highlighting}
\end{Shaded}

\begin{itemize}
\item
  Solo se van a permitir 3 hilos simultáneos.
\item
  Se crean 10 hilos y se guardan en un vector los objetos
  \{cpp\_thread\} para luego poder esperar por ellos usando
  \{cpp\_thread\_join\}, antes de terminar el programa.
\item
  Se llama a \texttt{semaphore::acquire()} antes de entrar en la
  \textbf{sección crítica}. Si hay más de 3 hilos en ella, el hilo es
  suspendido en la \textbf{variable de condición} \texttt{cv\_} de
  \texttt{semaphore}.
\item
  Se llama a \texttt{semaphore::release()} al salir de la
  \textbf{sección critica}. Esto hace que la \textbf{variable de
  condición} \texttt{cv\_} de \texttt{semaphore} despierte a uno de los
  hilos que está esperando en \texttt{semaphore::acquire()} para que
  entre en la \textbf{sección crítica}.
\end{itemize}

\section{Esperas}\label{_esperas}

Muchos de los objetos de sincronización que hemos visto necesitan algún
mecanismo para poner en espera a los hilos que los usan. Existen dos
alternativas desde el punto de vista de cómo implementar esta espera:

\begin{itemize}
\item
  El sistema operativo puede cambiar el estado del hilo o proceso a
  \textbf{esperado} y moverlo a una cola de espera asociada al objeto de
  sincronización, tal y como hemos comentado en varias ocasiones.
  Entonces el planificador de la CPU escogerá a otro proceso para ser
  ejecutado.
\item
  El hilo puede iterar comprobando constantemente la condición hasta que
  se cumple. A esa técnica se la denomina \textbf{espera ocupada}{}.
\end{itemize}

Este tipo de \textbf{espera ocupada} desperdicia tiempo de CPU que otro
hilo podría utilizar de forma más productiva, por lo que solo se utiliza
en el caso de esperas previsiblemente cortas. Para evitar que las
esperas ocupadas sean demasiado largas, los sistemas operativos nunca
expulsan de la CPU a hilos que se estén ejecutando dentro de secciones
críticas controladas por objetos de sincronización con este tipo de
espera, con la idea de que salgan de la sección crítica lo antes
posible.

En Windows API, por ejemplo, se puede utilizar
\{win32\_initializecriticalsectionandspincount\} para inicializar un
objeto de \textbf{sección crítica} donde el hilo que la intenta adquirir
itera el número especificado de veces en una \textbf{espera ocupada},
comprobando si la sección es liberada, antes de bloquearse en el estado
\textbf{esperando} si eso no ocurre.

La \textbf{espera ocupada} de estos objetos \textbf{sección crítica} de
Windows API solo ocurre en sistemas multiprocesador, donde el hilo que
tiene adquirida la sección puede estar ejecutándose en otro hilo y
terminar rápidamente. En sistemas monoprocesador nunca hay
\textbf{espera ocupada}.

A los \textbf{\emph{mutex}} con \textbf{espera ocupada} también se los
denomina \textbf{\emph{{}spinlock}}. Los \textbf{\emph{spinlocks}} son
utilizados frecuentemente para proteger las estructuras del núcleo en
los sistemas multiprocesador, cuando la tarea a realizar dentro de la
sección crítica en el núcleo requiere poco tiempo y los diseñadores
calculan que se desperdicia más tiempo sacando de la CPU al hilo en
espera para ejecutar otro en su lugar.

\section{Funciones reentrantes y seguras en
hilos}\label{_funciones_reentrantes_y_seguras_en_hilos}

Todas estas cuestiones sobre la sincronización no solo afectan al código
que escribimos sino también a las librerías que podemos utilizar. A la
hora de decidir utilizar una librería en un programa multihilo es
necesario que tengamos en cuenta los conceptos de \textbf{reentrante} y
\textbf{seguridad de hilos}.

\subsection{Funciones reentrantes}\label{_funciones_reentrantes}

{} {} Una función es \textbf{reentrante} puede ser interrumpida en medio
de su ejecución y, mientras espera, volver a ser llamada con total
seguridad. Obviamente las funciones recursivas deben ser reentrantes
para poder llamarse a sí mismas una y otra vez con seguridad.

En el contexto de la programación multihilo, ocurre una reentrada cuando
durante la ejecución de una función por parte de un hilo, este es
interrumpido por el sistema operativo para planificar posteriormente a
otro del mismo proceso que invoca la misma función.

En general una función es reentrante, si:

\begin{itemize}
\item
  No modifica variables estáticas o globales. Si lo hiciera solo puede
  hacerlo mediante operaciones \textbf{leer-modificar-escribir} que sean
  ininterrumpibles ---es decir, atómicas---.
\item
  No modifica su propio código y no llama a otras funciones que no sean
  reentrantes.
\end{itemize}

Como hemos mencionado anteriormente, los \textbf{manejadores de señal}
deben ser funciones \textbf{reentrantes}.

\subsection{Seguridad en hilos}\label{_seguridad_en_hilos}

{} {} Una función es \textbf{segura en hilos} o \textbf{{}thread-safe}
si al manipular estructuras compartidas de datos lo hace de tal manera
que se garantiza la ejecución segura de la misma por múltiples hilos al
mismo tiempo. Obviamente estamos hablando de un problema de secciones
críticas, por lo que las funciones lo resuelven sincronizando el acceso
a estos datos mediante el uso de \textbf{semáforos},
\textbf{\emph{mutex}} u otros recursos similares ofrecidos por el
sistema operativo.

En ocasiones, ambos conceptos se confunden porque es bastante común que
el código reentrante también sea seguro en hilos. Sin embargo es posible
crear código reentrante que no sea seguro en hilos y viceversa.

A la hora de usar una función o librería que va a ser llamada desde
múltiples hilos, primero debemos consultar la documentación para
averiguar si es \textbf{segura en hilos}. Si no lo fuera, tendríamos que
buscar funciones alternativas o recordar proteger las llamadas a las
funciones no seguras con mecanismos de sincronización, para asegurar que
solo son invocadas desde un hilo al mismo tiempo.

Esto se aplica tanto a librerías de otros desarrolladores como a la
librería estándar del lenguaje que estemos usando y a la librería del
sistema.

\subsubsection{Seguridad en hilos en
C++}\label{_seguridad_en_hilos_en_c}

La norma general es que las clases de la librería estándar de C++ son
seguras frente a múltiples accesos de lectura desde diferentes hilos.
Pero si un hilo modifica un objeto, todas las lecturas y escrituras en
el mismo objeto por ese y otros hilos deben estar protegidas.

Obviamente, las clases de mecanismos de sincronización y gestión de
hilos generalmente ofrecen mayores garantías, para lo que hay que
consultar la documentación.

\subsubsection{Seguridad en hilos en
C}\label{_seguridad_en_hilos_en_c_2}

El estándar de C no menciona nada sobre \textbf{seguridad en hilos}, por
lo que se debe suponer que las funciones de la librería estándar no lo
son o consultar la documentación ofrecida por el proveedor de la
librería.

Por ejemplo, todas las versiones de la librería estándar de C en Windows
actualmente son \textbf{seguras en
hilos}\hyperref[Microsoft-MultithreadingWithC]{???}. Pero hasta hace
unos años Microsoft ofrecía varias versiones de la librería, algunas
\textbf{seguras en hilos}, para usar en aplicaciones multihilo, y otras
no seguras para usar en aplicaciones monohilo.

\subsubsection{Seguridad en hilos en
POSIX}\label{_seguridad_en_hilos_en_posix}

La API POSIX es un superconjunto de la API de la librería estándar de C.
Por lo que en esos sistemas el estándar POSIX es el que marca qué
funciones de la librería estándar de C y del resto de la API POSIX son
\textbf{seguras en hilos}.

En los sistemas POSIX, el estándar establece que todas las funciones son
seguras excepto algunas muy concretas, que se pueden consultar en el
apartado
\href{https://pubs.opengroup.org/onlinepubs/9699919799/functions/V2_chap02.html\#tag_15_09_01}{«2.9.1
Thread-Safety»} de la especificación. Muchas de esas funciones no se
especifican como \textbf{seguras en hilos} porque existe alguna
alternativa que sí lo es. Por ejemplo, \{linux\_strerror\} no es
\textbf{segura en hilos}, pero \{linux\_strerror\_r\} tiene una
funcionalidad equivalente y sí lo es.
